% =====================================================================
\section{Introduction}\label{sec:intro}
% =====================================================================
Taiwan's ``garbage never touches the ground'' policy requires residents
to hand their garbage directly to a collection truck during its scheduled
stop. In Taichung City, the official real-time service has run for more
than ten years. As of September 2026 it covers 29 districts, more than
370 routes, and over 20{,}000 collection points, and its app has been
downloaded nearly 650{,}000 times~\cite{tcepb2026cleaner}. Despite this
demand, users are not satisfied with the timing information. The official
app has an average App Store rating of about 2.6 stars (609 ratings). Of
the 100 most recent reviews collected on 2026-11-07, 71 gave one star,
and 39 of those concerned inaccurate times or wrong arrival
information~\cite{appstore_official}.

A wrong decision is costly because the window is short. During the 2026
Lunar New Year holidays, the city's 701 fixed-time collection stations
each kept a truck for only about 5 to 30 minutes~\cite{tcepb2026spring}.
On national holidays, curbside collection stops and only fixed-time
points are served~\cite{tcepb_holiday}. A resident who misses the stop
must wait for the next round or find another point.

Two open datasets support third-party services: a vehicle-position feed
declared to update every 10 minutes~\cite{dataset83558} and a list of
collection points with scheduled times~\cite{dataset84004}. Existing
apps, however, answer only ``where is the truck'' or ``when will it
arrive at this point.'' The official app filters nearby points by a fixed
radius and shows scheduled times, without walking time. A popular private
app sorts nearby points by distance and shows walking time, but it does
not combine walking time with the truck's arrival, and its reminders use
a fixed lead time~\cite{appstore_private}. Neither app shows uncertainty
or suggests an alternative when a user cannot make it.

This paper addresses a decision question rather than a lookup:
\emph{Which collection point should the user head to now to be most
likely to catch the truck?} The nearest point is not always the best,
because the truck may be about to leave it while a farther point still
has time to spare. The answer depends on the truck's estimated time of
arrival (ETA), the user's walking time, and a safety buffer for data
delay and prediction error.

The main contributions of this paper are as follows:
\begin{itemize}
  \item We formulate ``catching the garbage truck'' as a time-feasibility
        recommendation problem that combines truck ETA, walking time, and
        uncertainty into a single slack measure.
  \item We design GarbageGo, a domain AI agent in which a large language
        model (LLM) interprets requests and explains results, while GIS,
        ETA, and recommendation engines compute every number.
  \item We report a data feasibility study of Taichung's open vehicle and
        collection-point data, covering update behavior, staleness,
        entity matching, and geocoding.
  \item We define an evaluation plan that compares nearest-point and
        slack-based strategies using catch rate and misjudgment rate.
\end{itemize}

The remainder of this paper is organized as follows.
Section~\ref{sec:related} reviews related work.
Section~\ref{sec:scheme} presents the proposed scheme.
Section~\ref{sec:exp} reports the data feasibility study and the
evaluation plan, and Section~\ref{sec:conclusion} concludes the paper.
